\subsection{Products of two quadratic character fields}
\label{sec:quadratic-product-bound}

The preceding moment bound also treats both quadratic factors at once.
The fields in the next theorem are allowed to coincide.

\begin{lemma}\label{lem:quadratic-codegree-three}
Let $R$ be a self-dual fusion ring. A character whose value field is
quadratic cannot have formal codegree three if its conjugate also has
formal codegree three.
\end{lemma}
\begin{proof}
First, a real quadratic algebraic integer $\alpha$ for which both
conjugates have absolute value at most $\sqrt2$ belongs to
\[
 \{0,1,-1,\sqrt2,-\sqrt2\}.
\]
For a rational $\alpha$ this is immediate. Otherwise write its minimal
polynomial as $t^2-Tt+N$. Its trace satisfies $|T|\leq2$.
If $T=0$, then $-N$ is a positive nonsquare integer at most two, so
$\alpha=\pm\sqrt2$. If $|T|=1$, irreducibility and real roots force
$N\leq-1$; one root then has absolute value at least
$(1+\sqrt5)/2>\sqrt2$. If $|T|=2$, they force $N\leq0$, and one
root has absolute value at least two. These contradictions prove the
claim.

Suppose $\chi$ and its quadratic conjugate $\chi'$ both have codegree
three. Since the ring is self-dual, the codegree is the sum of the
squares of the real character values. The unit contributes one, so
each nonunit value and its conjugate have absolute value at most
$\sqrt2$. Some value is irrational because the character's value
field is quadratic. The preceding claim makes it $\pm\sqrt2$,
and the remaining nonunit values must vanish. Thus the two character
vectors, in a suitable common ordering of the basis, are
\[
 (1,\sqrt2,0,\ldots,0),\qquad
 (1,-\sqrt2,0,\ldots,0).
\]
Their scalar product is $-1$, contradicting the orthogonality of
distinct characters of a self-dual fusion ring.
\end{proof}

\begin{lemma}\label{lem:small-quadratic-codegree-norm}
Let $E,E'$ be the formal codegrees of a quadratic character orbit of a
complex-categorifiable self-dual rank-four fusion ring. Then
$EE'\geq16$.
\end{lemma}
\begin{proof}
Set $m=EE'$ and $B=E+E'$. Both are positive integers. The formal
codegrees are totally positive, so $B^2\geq4m$. Their reciprocal sum
is strictly less than one, since there are two other characters and
all codegrees are positive. Hence $B<m$.

The $d$-number criterion gives $m\mid B^2$. If $E$ is nonrational,
apply \cite[Proposition~\ref{B-prop:standard}]{Finite} to $t^2-Bt+m$. If $E$ is
rational, then $E=E'$ is an integer and the same divisibility follows
from $B=2E$ and $m=E^2$.

Here is a short exact small-norm argument. Write $m=ab^2$, where $a$
is a positive squarefree integer and $b$ is a positive integer. The
divisibility $m\mid B^2$ implies $B=ab\,t$ for a positive integer
$t$. The inequalities above become
\[
 t<b,\qquad at^2\geq4.
\]
If $m<16$, then $b\leq3$. The case $b=1$ is impossible. If $b=2$,
then $a\leq3$ and $t=1$, contradicting $at^2\geq4$. If $b=3$,
then $a=1$ and the only possibility is $t=2$. Thus the unique
remaining pair is $(m,B)=(9,6)$, which gives $E=E'=3$.
Lemma~\ref{lem:quadratic-codegree-three} excludes this pair.
\end{proof}

\begin{theorem}\label{thm:quadratic-product-bound}
Let $R$ be a self-dual rank-four fusion ring with
\[
 R\otimes_{\mathbb Z}\mathbb Q\cong K_1\times K_2,
\]
where $K_1,K_2$ are real quadratic fields, possibly equal. If $R$
admits a pivotal categorification, then
\[
 \operatorname{FPdim}(R)<196,\qquad \mu(R)\leq13.
\]
The only pivotal-categorifiable rings of this algebra type are
$\mathrm{Fib}\boxtimes\operatorname{Vec}_{\mathbb Z/2}$ and
$\mathrm{Fib}\boxtimes\mathrm{Fib}$, and both admit unitary
categorifications.
\end{theorem}
\begin{proof}
By Corollary~\ref{cor:pivotal-pu-unrestricted}, a pseudo-unitary
categorification exists. Number the characters so that $D,D'$ are
the codegrees of the Frobenius--Perron orbit and $E,E'$ are the
codegrees of the other quadratic orbit. Conjugate codegrees are
allowed to coincide. Put
\[
 m=EE',\qquad B=E+E',\qquad u=B/m,\qquad c=1-u.
\]
The reciprocal-codegree identity gives $1/D'=c-1/D$, with $0<c<1$.
Consequently the quantity in Lemma~\ref{lem:codegree-moment-bound}
satisfies
\begin{align*}
 w&=\frac1{D^2}+\frac1{(D')^2}+\frac1{E^2}+\frac1{(E')^2}\\
  &=u^2+(1-u)^2-\frac2m-\frac{2c}{D}+\frac2{D^2}\\
  &\geq\frac12-\frac2m-\frac2D
   \geq\frac38-\frac2D,
\end{align*}
where the last inequality is
Lemma~\ref{lem:small-quadratic-codegree-norm}.

Corollary~\ref{cor:dimension-196} gives $D<196$ and multiplicity at most thirteen.

The bounded-classification certificates through multiplicity
thirteen now apply. Of the positive self-dual rings, precisely
$F_{1,2}$ and $F_{1,5}$ have a rational algebra consisting of two
quadratic factors. They are the two rings in the statement. For the
first, the algebra is the tensor product of the Fibonacci quadratic
field with $\mathbb Q[\mathbb Z/2]\cong\mathbb Q\times\mathbb Q$;
for the second, it is
$\mathbb Q(\sqrt5)\otimes_{\mathbb Q}\mathbb Q(\sqrt5)$.
Both are $\mathbb Q(\sqrt5)\times\mathbb Q(\sqrt5)$ and have the
unitary models given in the realization table.
\end{proof}

\begin{corollary}\label{cor:distinct-quadratics-excluded}
No self-dual rank-four fusion ring with rational algebra
$K_1\times K_2$, for distinct real quadratic fields $K_1,K_2$,
admits a pivotal categorification. No multiplicity restriction is
required.
\end{corollary}
